\section{Evaluation Setup}
\label{sec:setup}
\label{sec:problem}

Given a passage $x$ and ordered candidates $C=(c_1,\ldots,c_m)$,
admission stores $K\subseteq C$ with $|K|\leq k$. The future question
$q$ and answer $a$ are hidden. When $q$ arrives, the frozen reader sees
only $q$ and $K$; the original passage has left its window. Every policy
receives the same candidates and budget. We measure three successive
outcomes: candidate ordering, retained information and answer correctness.

\paragraph{Ranking within the episode.}
Within-episode area under the receiver operating characteristic curve
(AUC) is the fraction of useful--distractor pairs ordered correctly,
with half credit for ties, averaged across episodes. For useful concept
$b_i$ in episode $i$, score $s$ and $n$ episodes,
\begin{equation}
A_{\mathrm{episode}}
=\frac{1}{n}\sum_{i=1}^{n}\frac{1}{|C_i|-1}
\sum_{c\in C_i\setminus\{b_i\}}
\left[\mathbf{1}\{s(b_i)>s(c)\}
+\tfrac12\mathbf{1}\{s(b_i)=s(c)\}\right].
\label{eq:episodeauc}
\end{equation}
It measures the ordering among concepts that actually compete for memory. Pooled AUC
instead compares candidates across the entire battery and summarizes
score discrimination. Retention at budget $k$ connects this ordering
to the admitted set. For the provenance analysis,
we replace utility labels with statedness labels and report
\emph{provenance AUC}: discrimination of stated from unstated candidates
under a case-insensitive stem match against the passage.

\paragraph{Retention of the useful concept.}
Bridge retention is the fraction of episodes whose stored set contains
the construction-labelled useful concept:
\begin{equation}
B_k=\frac{1}{n}\sum_{i=1}^{n}\mathbf{1}\{b_i\in K_i\}.
\label{eq:retention}
\end{equation}
It directly measures what
survives the budget; correct answering additionally depends on the
reader and the other stored concepts. An oracle-label policy prioritizes
the intended bridge and provides a reference for this construction.

\paragraph{Downstream recall.}
Recall is the fraction of future questions answered correctly from
stored concepts by the frozen reader. It measures the downstream effect
of admission for a fixed reader. A policy can rank the bridge
above most distractors yet drop it at a small budget; a retained bridge
can still lead to an incorrect answer. Reporting all three outcomes
locates these distinct transitions from score to memory to answer.
The reader answers greedily, with Qwen3 thinking disabled. A blind
semantic judge receives only the question, gold answer and response;
it accepts a committed correct answer or exact equivalent. A separate
string-match grader checks agreement (App.~\ref{app:judge}).
Question-only and whole-passage readers establish how much the answer
can be recovered before admission and with the complete original input.

\paragraph{Comparisons.}
Baselines are the untrained yes/no importance report, a 1--10 importance
rating, reversed yes/no report, admit what the passage omits,
omission then report, wording variants, the provenance-adjusted report,
the combined teacher used directly, random selection, output likelihood,
the workspace readout, and students trained from either teacher signal
alone. All receive the passage with the future question hidden
(App.~\ref{app:baselines}). Paired AUC differences use 4,000-draw
whole-episode bootstraps and 95\% intervals; recall comparisons use
exact McNemar tests on paired answers and seed-averaged differences
with paired 95\% bootstrap intervals (Table~\ref{tab:paired}). Trained results give three-seed
means and sample standard deviations. Bold table entries identify SDA.

\paragraph{Evaluation scale.}
\label{sec:breadth}
\label{sec:models}

The principal diagnostic comparison covers eight models from four
families: Qwen2.5-7B, OLMo-2-7B, Mistral-7B and Qwen3 at 1.7B, 4B,
8B, 14B and 32B
\citep{qwen2025,olmo2,jiang2023mistral,qwen3}. The extended grid contains
23 checkpoints across those families and GPT-2, including base and
post-training stages (App.~\ref{app:breadth}). SDA trains three seeds
for each of six models: Qwen3-8B, 14B and 32B and the other three
7B families, totaling 18 principal training runs. Each supplies its
own teacher, adapter and frozen reader.

Training and validation contain 175 and 45 episodes, drawn from
88 Explicit episodes and 132 implied-answer episodes, including the
57-episode alternate-author benchmark. In an implied-answer episode,
the passage cues an unstated useful concept that itself supplies the
answer; Decoupled instead requires a further inference from that concept. Development uses Decoupled
(68 episodes), Compositional (52), and two separate 68-episode pools
that add confusable or unrelated unstated distractors to Decoupled.
Fresh-1 and Fresh-2 contain 316 and 363 episodes; the entity-bridge
benchmark contains 42. Fresh-2 passage and pool controls reuse its
363 episodes. Budgets are $k=1,2,3$; the main comparison uses the
statedness-balanced pool, with $k=2$ reported across all six models.
The paired Fresh-2 grading audit totals 13,068 answers across the
18 SDA--report comparisons; both graders evaluate every answer.


